\documentclass[aspectratio=169]{beamer}
\usepackage[english]{babel}
\usepackage[utf8]{inputenc}
\usepackage[T1]{fontenc}
\usepackage{lmodern}
\usepackage{listings}
\usepackage{xcolor}
\usepackage{graphicx}
\usepackage{textcomp}
\DeclareUnicodeCharacter{2014}{---}
\DeclareUnicodeCharacter{2013}{--}
\DeclareUnicodeCharacter{2212}{-}
\DeclareUnicodeCharacter{2192}{\ensuremath{\rightarrow}}
\DeclareUnicodeCharacter{00B1}{\ensuremath{\pm}}
\DeclareUnicodeCharacter{00D7}{\ensuremath{\times}}
\DeclareUnicodeCharacter{00B7}{\ensuremath{\cdot}}
\DeclareUnicodeCharacter{20AC}{EUR}
\DeclareUnicodeCharacter{2070}{\textsuperscript{0}}
\DeclareUnicodeCharacter{00B9}{\textsuperscript{1}}
\DeclareUnicodeCharacter{00B2}{\textsuperscript{2}}
\DeclareUnicodeCharacter{00B3}{\textsuperscript{3}}
\DeclareUnicodeCharacter{2074}{\textsuperscript{4}}
\DeclareUnicodeCharacter{2075}{\textsuperscript{5}}
\DeclareUnicodeCharacter{2076}{\textsuperscript{6}}
\DeclareUnicodeCharacter{2077}{\textsuperscript{7}}
\DeclareUnicodeCharacter{2078}{\textsuperscript{8}}
\DeclareUnicodeCharacter{2079}{\textsuperscript{9}}
\DeclareUnicodeCharacter{2248}{\ensuremath{\approx}}
\DeclareUnicodeCharacter{2264}{\ensuremath{\le}}

% Colors for code
\definecolor{codebg}{RGB}{245,245,245}
\definecolor{codecomment}{RGB}{0,128,0}
\definecolor{codekeyword}{RGB}{0,0,255}
\definecolor{codestring}{RGB}{163,21,21}
\definecolor{bandcolor}{RGB}{200,50,50}

\lstdefinestyle{pythonstyle}{
    language=Python,
    backgroundcolor=\color{codebg},
    basicstyle=\ttfamily\scriptsize,
    keywordstyle=\color{codekeyword},
    commentstyle=\color{codecomment},
    stringstyle=\color{codestring},
    showstringspaces=false,
    breaklines=true,
    frame=single,
    rulecolor=\color{black},
    escapeinside={(*@}{@*)},
    moredelim=**[l][\color{bandcolor}\bfseries]{\#\ ---\ NEW},
    moredelim=**[l][\color{bandcolor}\bfseries]{\#\ ---\ CHANGED},
    literate=
      {á}{{\'a}}1 {é}{{\'e}}1 {í}{{\'i}}1 {ó}{{\'o}}1 {ú}{{\'u}}1
      {Á}{{\'A}}1 {É}{{\'E}}1 {Í}{{\'I}}1 {Ó}{{\'O}}1 {Ú}{{\'U}}1
      {ñ}{{\~n}}1 {Ñ}{{\~N}}1 {ü}{{\"u}}1 {Ü}{{\"U}}1
      {¿}{{?`}}1 {¡}{{!`}}1
      {—}{{---}}1 {–}{{--}}1 {−}{{-}}1
      {±}{{\ensuremath{\pm}}}1 {×}{{\ensuremath{\times}}}1
      {→}{{\ensuremath{\rightarrow}}}1
      {°}{{\textdegree}}1
      {·}{{\ensuremath{\cdot}}}1
      {€}{{\texteuro}}1
      {≈}{{\ensuremath{\approx}}}1
      {≤}{{\ensuremath{\le}}}1
      {⁰}{{\textsuperscript{0}}}1 {¹}{{\textsuperscript{1}}}1
      {²}{{\textsuperscript{2}}}1 {³}{{\textsuperscript{3}}}1
      {⁴}{{\textsuperscript{4}}}1 {⁵}{{\textsuperscript{5}}}1
      {⁶}{{\textsuperscript{6}}}1 {⁷}{{\textsuperscript{7}}}1
      {⁸}{{\textsuperscript{8}}}1 {⁹}{{\textsuperscript{9}}}1
}

\lstset{style=pythonstyle}

% Title slide macro: \lessontitle{Lesson Number}{Title}
\newcommand{\lessontitle}[2]{
    \title{Lesson #1: #2}
    \author{}
    \institute{Tec de Monterrey}
    \begin{frame}[plain]
        \titlepage
    \end{frame}
}

% Figure frame macro: \figureframe{Title}{Image path}
\newcommand{\figureframe}[2]{
    \begin{frame}{#1}
        \begin{center}
            \includegraphics[height=0.8\textheight,width=\textwidth,keepaspectratio]{#2}
        \end{center}
    \end{frame}
}


% Listings pulled from src/ with \lstinputlisting, so the slide is the file.
\lstset{
    basicstyle=\ttfamily\fontsize{8pt}{9.6pt}\selectfont,
    breaklines=true,
    breakatwhitespace=true,
    columns=fullflexible,
    keepspaces=true,
    xleftmargin=0.3em,
    framesep=2pt,
    aboveskip=0pt,
    belowskip=0pt
}

\newcommand{\resultframe}[3]{%
    \begin{frame}{#1}
        \begin{center}
            \includegraphics[height=0.62\textheight,width=\textwidth,keepaspectratio]{#2}
        \end{center}
        \vspace{0.25em}
        {\small #3\par}
    \end{frame}%
}

\newcommand{\claimframe}[2]{%
    \begin{frame}{#1}
        {\small #2\par}
    \end{frame}%
}

\date{}

\begin{document}

\lessontitle{03}{Selection}

\section{This lesson}

\begin{frame}{This lesson}
Lesson 02 named the five phases. Selection is the first operator to look at on its own, still without measuring a whole run.

\vspace{0.6em}
\begin{itemize}
    \item Proportional, rank, SUS, tournament: different ways to discard.
    \item Elitism is the cheap insurance that the best so far survives.
    \item One generation is not a verdict; that measurement is Lesson 06.
\end{itemize}
\end{frame}

% ================= selection_pressure_01_the_population =================
\begin{frame}{Selection probabilities}
\[
p_i=\frac{w_i}{\sum_j w_j},\qquad \mathbb E[N_i]=Np_i,
\qquad \Pr(R=r)=\left(\frac{N-r+1}{N}\right)^k-
\left(\frac{N-r}{N}\right)^k .
\]
Proportional selection changes when its shift floor changes. Rank selection
uses order instead of scale; SUS changes sampling variance, not expectation.
Here (r=1) is best; tournament probabilities telescope to one. Elitism preserves the best only for stationary deterministic fitness and an
unchanged elite copy.
\end{frame}

\section{The population, and how to measure what selection does}

\begin{frame}{The population, and how to measure what selection does}
Selection never invents anything. It takes a population of N and returns a
population of N, built only from copies of what was already there. So the only
thing worth measuring is the BOOKKEEPING: who got copied, how many times, and
who disappeared.
\end{frame}

\resultframe{The population, and how to measure what selection does}{../figures/selection_pressure_01_the_population.png}{%
    Copying everyone once: 0 extinct, diversity 6.665, 1 copy of the best — the zero of every later measurement. Fitness is negative almost everywhere, so no wheel can be built yet}

% ================= selection_pressure_02_proportional =================
\section{Proportional selection, and the floor nobody mentions}

\begin{frame}{Proportional selection, and the floor nobody mentions}
The standard picture of selection is a roulette wheel whose sectors are as wide
as the individuals' fitness. On this population the picture does not fit: every
fitness is negative, and a sector cannot have a negative width. The usual
repair is to shift all the values up until the worst one is non-negative - and
that repair silently introduces a parameter, the floor, which turns out to
control the selection pressure more strongly than the fitness values do.
\end{frame}

\begin{frame}[fragile]{(1) shift\_to\_positive()}
\lstinputlisting[basicstyle=\ttfamily\fontsize{8pt}{9.6pt}\selectfont,firstline=119,lastline=130]{../src/selection_pressure_02_proportional.py}
\end{frame}

\begin{frame}[fragile]{(2) select\_proportional()}
\lstinputlisting[basicstyle=\ttfamily\fontsize{7pt}{8.5pt}\selectfont,firstline=133,lastline=154]{../src/selection_pressure_02_proportional.py}
\end{frame}

\begin{frame}[fragile]{(3) pressure\_report()}
\lstinputlisting[basicstyle=\ttfamily\fontsize{6pt}{7.2pt}\selectfont,firstline=157,lastline=205]{../src/selection_pressure_02_proportional.py}
\end{frame}

\begin{frame}[fragile]{(4) the two floors}
\lstinputlisting[basicstyle=\ttfamily\fontsize{6pt}{7.2pt}\selectfont,firstline=208,lastline=260]{../src/selection_pressure_02_proportional.py}
\end{frame}

\begin{frame}[fragile]{(5) the wheel figure}
\lstinputlisting[basicstyle=\ttfamily\fontsize{7pt}{8.5pt}\selectfont,firstline=262,lastline=285]{../src/selection_pressure_02_proportional.py}
\end{frame}

\resultframe{Proportional selection, and the floor nobody mentions}{../figures/selection_pressure_02_proportional.png}{%
    The floor *is* the pressure: the best expects \textbf{2.35} copies at floor 0.001 but \textbf{1.30} at floor 3.0 — same population, same wheel}

% ================= selection_pressure_03_rank =================
\section{Rank selection, or pressure that does not depend on scale}

\begin{frame}{Rank selection, or pressure that does not depend on scale}
Step 2 left an embarrassment: the same wheel, on the same population, applies
very different pressure depending on a floor we invented. Rank selection throws
the fitness values away and keeps only their ORDER. The sector widths then
depend on the population size alone, so the pressure is the same whatever the
fitness scale is - and in particular it cannot be changed by adding a constant.
\end{frame}

\begin{frame}[fragile]{(1) select\_rank()}
\lstinputlisting[basicstyle=\ttfamily\fontsize{7pt}{8.5pt}\selectfont,firstline=199,lastline=221]{../src/selection_pressure_03_rank.py}
\end{frame}

\begin{frame}[fragile]{(2) the rank row}
\lstinputlisting[basicstyle=\ttfamily\fontsize{6pt}{7.2pt}\selectfont,firstline=227,lastline=279]{../src/selection_pressure_03_rank.py}
\end{frame}

\begin{frame}[fragile]{(3) the sector figure}
\lstinputlisting[basicstyle=\ttfamily\fontsize{6pt}{7.2pt}\selectfont,firstline=281,lastline=312]{../src/selection_pressure_03_rank.py}
\end{frame}

\resultframe{Rank selection, or pressure that does not depend on scale}{../figures/selection_pressure_03_rank.png}{%
    Order alone fixes the pressure at one definite value, immune to rescaling — but the best individual is still lost in \textbf{13.6\%} of draws}

% ================= selection_pressure_04_elitism =================
\section{Elitism, the cheapest guarantee in the course}

\begin{frame}{Elitism, the cheapest guarantee in the course}
Steps 2 and 3 both share a flaw that is easy to overlook: the best individual
is drawn at random like everybody else, so it is sometimes not drawn at all. A
generation can be strictly worse than its parent generation. Elitism fixes that
by copying the top few individuals into the new generation before any draw
happens. It is one line of code, it removes the regression completely, and it
costs diversity - all three of which this step measures.
\end{frame}

\begin{frame}[fragile]{(1) select\_rank\_with\_elite()}
\lstinputlisting[basicstyle=\ttfamily\fontsize{7pt}{8.5pt}\selectfont,firstline=224,lastline=247]{../src/selection_pressure_04_elitism.py}
\end{frame}

\begin{frame}[fragile]{(2) the elite rows}
\lstinputlisting[basicstyle=\ttfamily\fontsize{6pt}{7.2pt}\selectfont,firstline=253,lastline=307]{../src/selection_pressure_04_elitism.py}
\end{frame}

\begin{frame}[fragile]{(3) the elitism figure}
\lstinputlisting[basicstyle=\ttfamily\fontsize{6pt}{7.2pt}\selectfont,firstline=309,lastline=337]{../src/selection_pressure_04_elitism.py}
\end{frame}

\resultframe{Elitism, the cheapest guarantee in the course}{../figures/selection_pressure_04_elitism.png}{%
    One line takes the loss of the best from 13.6\% to \textbf{0.0\%}, raises pressure (1.80 → 2.64 copies) and costs diversity (5.448 → 5.210)}

% ================= selection_pressure_05_sus =================
\section{Stochastic universal sampling, or the same wheel spun once}

\begin{frame}{Stochastic universal sampling, or the same wheel spun once}
Proportional selection spins the wheel N times, so luck accumulates: an
individual owning a quarter of the wheel can still end up with nothing, and
another can win four slots in a row. SUS keeps exactly the same wheel and the
same expected number of copies, but takes all N samples in one go, with evenly
spaced pointers. The expectation does not move. The spread collapses.
\end{frame}

\begin{frame}[fragile]{(1) select\_stochastic\_universal\_sampling()}
\lstinputlisting[basicstyle=\ttfamily\fontsize{7pt}{8.5pt}\selectfont,firstline=247,lastline=271]{../src/selection_pressure_05_sus.py}
\end{frame}

\begin{frame}[fragile]{(2) the SUS rows}
\lstinputlisting[basicstyle=\ttfamily\fontsize{6pt}{7.2pt}\selectfont,firstline=277,lastline=339]{../src/selection_pressure_05_sus.py}
\end{frame}

\begin{frame}[fragile]{(3) the spread figure}
\lstinputlisting[basicstyle=\ttfamily\fontsize{6pt}{7.2pt}\selectfont,firstline=341,lastline=389]{../src/selection_pressure_05_sus.py}
\end{frame}

\resultframe{Stochastic universal sampling, or the same wheel spun once}{../figures/selection_pressure_05_sus.png}{%
    Same wheel, same expectation, spread collapses. And it does \textbf{not} fix the floor: 2.35 vs 1.31 copies, the same split step 2 found}

% ================= selection_pressure_06_tournament =================
\section{Tournament selection, and all five methods on one scale}

\begin{frame}{Tournament selection, and all five methods on one scale}
Tournament selection needs no wheel, no sort, no sum and no floor: pick a few
individuals at random, keep the fittest, repeat. This is the method Lesson 01
used, and with the same seed this file reproduces Lesson 01's draw exactly. Its
real virtue is the one thing none of the other four offer: a single integer that
sets the selection pressure, from none at all to almost total.
\end{frame}

\begin{frame}[fragile]{(1) select\_tournament()}
\lstinputlisting[basicstyle=\ttfamily\fontsize{8pt}{9.6pt}\selectfont,firstline=279,lastline=294]{../src/selection_pressure_06_tournament.py}
\end{frame}

\begin{frame}[fragile]{(2) the draw from Lesson 01}
\lstinputlisting[basicstyle=\ttfamily\fontsize{8pt}{9.6pt}\selectfont,firstline=300,lastline=313]{../src/selection_pressure_06_tournament.py}
\end{frame}

\begin{frame}[fragile]{(3) the pressure dial}
\lstinputlisting[basicstyle=\ttfamily\fontsize{8pt}{9.6pt}\selectfont,firstline=315,lastline=332]{../src/selection_pressure_06_tournament.py}
\end{frame}

\begin{frame}[fragile]{(4) the whole comparison}
\lstinputlisting[basicstyle=\ttfamily\fontsize{6pt}{7.2pt}\selectfont,firstline=334,lastline=388]{../src/selection_pressure_06_tournament.py}
\end{frame}

\begin{frame}[fragile]{(5) the dial figure}
\lstinputlisting[basicstyle=\ttfamily\fontsize{6pt}{7.2pt}\selectfont,firstline=390,lastline=421]{../src/selection_pressure_06_tournament.py}
\end{frame}

\resultframe{Tournament selection, and all five methods on one scale}{../figures/selection_pressure_06_tournament.png}{%
    One integer sets the pressure: k=3 → 2.74 copies, k=5 → 4.13, k=10 → 6.52. Across all 13 configurations, pressure vs quality \textbf{+0.95}, pressure vs diversity \textbf{−0.92}}

\section{Next}

\begin{frame}{Next}
The two worst individuals were not selected zero times --- the table in the step has to be read, not a slogan above it.

\vspace{0.8em}
\textbf{Lesson 04 --- Crossover}

The second operator, still measured on a fixed pair of parents. Cut-and-swap copies; arithmetic computes. Those families do not swap.
\end{frame}
\end{document}
